% TEMPLATE-MILC-v3.tex - plantilla maestra para todas las guias MILC v3
% Ver scripts/generadores/build-guia-milc.py para la lista completa de
% claves esperadas. El script valida que no queden placeholders sin
% reemplazar antes de escribir el archivo final.

\documentclass[11pt,letterpaper]{article}
\usepackage[margin=1.75cm]{geometry}
\usepackage[table]{xcolor}
\usepackage{fontspec}
\usepackage[spanish]{babel}
\usepackage{tikz}
\usepackage[most]{tcolorbox}
\usepackage{tabularx}
\usepackage{array}
\usepackage{fancyhdr}
\usepackage{titlesec}
\usepackage{ragged2e}
\usepackage{enumitem}
\usepackage{microtype}

\setmainfont{Helvetica Neue}
\setsansfont{Helvetica Neue}
\setlength{\parindent}{0pt}
\setlength{\parskip}{6pt}
\hyphenpenalty=200
\exhyphenpenalty=200
\pretolerance=400
\tolerance=2000
\setlength{\emergencystretch}{1em}
\renewcommand{\arraystretch}{1.25}
\setlist{leftmargin=5mm,itemsep=1mm,topsep=1mm}
\newcolumntype{Y}{>{\RaggedRight\arraybackslash}X}

\definecolor{milcMagenta}{HTML}{E6007E}
\definecolor{milcVino}{HTML}{5A0038}
\definecolor{milcTurquesa}{HTML}{0093A5}
\definecolor{milcVerde}{HTML}{58A923}
\definecolor{milcMostaza}{HTML}{E5B400}
\definecolor{milcOcre}{HTML}{B86600}
\definecolor{milcNegro}{HTML}{241816}
\definecolor{milcGris}{HTML}{F7F7F4}
\definecolor{milcLinea}{HTML}{E8E2DD}
\definecolor{milcGradePrimary}{HTML}{ED7446}
\definecolor{milcGradeDark}{HTML}{B63E16}
\definecolor{milcGradeSoft}{HTML}{FFF0E8}
\definecolor{milcGradeAccent}{HTML}{CFE7FF}
\definecolor{milcGradeText}{HTML}{FFFFFF}
\color{milcNegro}

\pagestyle{fancy}
\fancyhf{}
\lhead{\small Tecnología e Informática · Grado 8}
\rhead{\small Guía 2 · MILC v3}
\cfoot{\small\thepage}
\renewcommand{\headrulewidth}{0pt}

\newtcolorbox{titlebox}[1]{
  enhanced,colback=#1,colframe=#1,arc=7mm,boxrule=0pt,
  left=7mm,right=7mm,top=3mm,bottom=3mm,
  before skip=8pt,after skip=8pt,
  fontupper=\sffamily\bfseries\Large\color{white}
}

\newtcolorbox{softbox}[3]{
  enhanced,colback=#2,colframe=#1,borderline west={4pt}{0pt}{#1},
  arc=2mm,boxrule=.4pt,left=4mm,right=4mm,top=3mm,bottom=3mm,
  before skip=6pt,after skip=8pt,title={#3},coltitle=white,
  colbacktitle=#1,fonttitle=\sffamily\bfseries,fontupper=\RaggedRight,
  attach boxed title to top left={xshift=3mm,yshift=-2mm},
  boxed title style={arc=2mm, boxrule=0pt}
}

\newtcolorbox{drawbox}[2]{
  enhanced,colback=white,colframe=#1,arc=4mm,boxrule=1pt,
  left=5mm,right=5mm,top=4mm,bottom=4mm,before skip=8pt,after skip=8pt,
  title={#2},coltitle=white,colbacktitle=#1,fonttitle=\sffamily\bfseries,
  fontupper=\RaggedRight,
  attach boxed title to top left={xshift=4mm,yshift=-2mm},
  boxed title style={arc=3mm, boxrule=0pt}
}

\newtcolorbox{infoband}[1]{
  enhanced,colback=#1!8,colframe=#1!50,arc=2mm,boxrule=.5pt,
  left=4mm,right=4mm,top=2mm,bottom=2mm,before skip=4pt,after skip=4pt,
  fontupper=\small\RaggedRight
}

\newcommand{\linead}[1]{\noindent\color{milcLinea}\rule{#1}{0.5pt}\color{milcNegro}}
\newcommand{\respuesta}[1]{\par\vspace{2mm}\foreach \n in {1,...,#1}{\noindent\color{milcLinea}\rule{\linewidth}{0.5pt}\par\vspace{5mm}}\color{milcNegro}}
\newcommand{\checkbox}{\textcolor{milcGradeDark}{\fbox{\phantom{x}}}\hspace{5pt}}

\newcommand{\phasecard}[4]{
\begin{tcolorbox}[enhanced,colback=#3,colframe=#2,arc=3mm,boxrule=.5pt,height=3.05cm,valign=center,halign=center,left=1.5mm,right=1.5mm,top=2mm,bottom=2mm]
{\sffamily\bfseries\fontsize{22}{24}\selectfont\textcolor{#2}{#1}}\par
{\sffamily\bfseries\small #4}
\end{tcolorbox}
}

\begin{document}
\thispagestyle{empty}
\begin{tikzpicture}[remember picture,overlay]
  \fill[white] (current page.south west) rectangle (current page.north east);
  \path[fill=milcGradePrimary,rounded corners=16mm]
    ([xshift=.55cm,yshift=-.55cm]current page.north west) rectangle
    ([xshift=-.55cm,yshift=3.65cm]current page.south east);

  \node[anchor=north west,text=milcGradeText] at ([xshift=2.0cm,yshift=-1.85cm]current page.north west)
    {\sffamily\bfseries\fontsize{14}{16}\selectfont GRADO};
  \node[anchor=north west,text=milcGradeText,opacity=.82] at ([xshift=2.0cm,yshift=-2.75cm]current page.north west)
    {\sffamily\bfseries\fontsize{9}{11}\selectfont SERIE GUÍAS · TECNOLOGÍA E INFORMÁTICA};

  \begin{scope}[shift={([xshift=-3.05cm,yshift=-2.55cm]current page.north east)}]
    \fill[white,opacity=.80] (-.62,-.52) rectangle (.62,.52);
    \draw[milcGradeText,line width=1pt] (-.62,-.52) rectangle (.62,.52);
    \fill[milcGradeDark] (-.42,-.38) rectangle (-.22,.06);
    \fill[milcGradeAccent] (-.10,-.38) rectangle (.10,.24);
    \fill[milcGradePrimary!60!black] (.22,-.38) rectangle (.42,.38);
  \end{scope}

  \node[anchor=west,text=milcGradeText] at ([xshift=1.9cm,yshift=-12.85cm]current page.north west)
    {\sffamily\bfseries\fontsize{124}{124}\selectfont 8};
  \draw[milcGradeText,line width=1.7pt]
    ([xshift=4.52cm,yshift=-11.68cm]current page.north west) circle (.13cm);

  \node[anchor=west,text=milcGradeText,text width=8.7cm,align=left] at ([xshift=10.35cm,yshift=-11.6cm]current page.north west)
    {
     {\sffamily\bfseries\fontsize{19}{22}\selectfont Tipos de datos y formato de celdas — ingreso y organización de información numérica}\\[4mm]
     {\sffamily\bfseries\fontsize{23}{26}\selectfont Guía 2-8-TIC}\\[2mm]
     {\sffamily\fontsize{13}{16}\selectfont Periodo Primero}\\[5mm]
     {\sffamily\bfseries\fontsize{11}{13}\selectfont Institución Educativa Sor María Juliana}\\[1mm]
     {\sffamily\fontsize{10}{12}\selectfont Autor: Dr. Álvaro Cárdenas Orozco}
    };

  \path[fill=milcGradeSoft,rounded corners=7mm]
    ([xshift=1.35cm,yshift=.85cm]current page.south west) rectangle
    ([xshift=-1.35cm,yshift=3.15cm]current page.south east);
  \draw[milcGradeDark,line width=.65pt,rounded corners=7mm]
    ([xshift=1.35cm,yshift=.85cm]current page.south west) rectangle
    ([xshift=-1.35cm,yshift=3.15cm]current page.south east);
  \node[anchor=center,text=milcNegro,text width=17.0cm,align=center] at ([yshift=2.0cm]current page.south)
    {\sffamily\fontsize{10}{12}\selectfont Metodología MILC · Apertura ancestral · Pensamiento computacional · Triángulo Dussel-Estoicismo-Floridi\\
    \textcolor{milcGradeDark}{\bfseries Producto:} Tabla en Excel con mínimo 4 columnas (texto, número, fecha, moneda) y 10 registros del entorno escolar correctamente formateados.};
\end{tikzpicture}
\mbox{}
\newpage

\section*{Apertura: Tipos de datos y formato de celdas — ingreso y organización de información numérica}

\begin{softbox}{milcGradeDark}{milcGradeSoft}{Saber ancestral}
El \textbf{cuaderno del tendero del barrio}: cada tipo de mercancía tiene su columna. Los plátanos por unidad, el arroz por libra, los huevos por docena, los cuadernos por número. Si confunde libras con unidades, pierde dinero al final del mes. La organización por tipo de medida es saber ancestral comercial.
\end{softbox}

\begin{softbox}{milcTurquesa}{milcGris}{Saber contemporáneo}
En \textbf{Excel y otras hojas de cálculo} cada tipo de dato (número, texto, fecha, moneda, porcentaje) tiene su formato específico. Aplicar mal el formato significa que el cálculo falla, los gráficos salen torcidos y la información engaña al que la mira.
\end{softbox}

\begin{softbox}{milcMostaza}{milcGris}{Pregunta-puente}
\textit{¿Qué del oficio del tendero — su disciplina de poner cada cosa en su columna — podemos llevar a la hoja de cálculo digital?}
\end{softbox}

\begin{softbox}{milcVerde}{milcGris}{Saber y saber hacer}
Al terminar podrás \textbf{distinguir tipos de datos} (número, texto, fecha, moneda), \textbf{aplicar el formato correcto} a cada celda, \textbf{nombrar bien los encabezados} de columna y \textbf{detectar errores} cuando un tipo y un formato no coinciden.
\end{softbox}

\small
\begin{tabularx}{\textwidth}{>{\bfseries}p{3.0cm}Y}
\rowcolor{milcGradePrimary!12}Autor & Dr. Álvaro Cárdenas Orozco\\
DBA & Estructura información numérica y categórica con tipos y formatos correctos en hojas de cálculo.\\
Referentes & MEN Tecnología · Estándares ISTE · Modelo MILC · Floridi (cuidado de datos) · ICONTEC ISO 9001 (registros)\\
Producto final & Tabla en Excel con mínimo 4 columnas (texto, número, fecha, moneda) y 10 registros del entorno escolar correctamente formateados.\\
\end{tabularx}
\normalsize

\begin{titlebox}{milcGradeDark}
Ruta MILC: así vamos a aprender
\end{titlebox}
\begin{tabularx}{\textwidth}{>{\centering\arraybackslash}X>{\centering\arraybackslash}X>{\centering\arraybackslash}X>{\centering\arraybackslash}X}
\phasecard{1}{milcVerde}{milcGris}{Escuta\\[-1mm]\small Observo, escucho y pregunto.} &
\phasecard{2}{milcTurquesa}{milcGris}{Sistematización\\[-1mm]\small Pensamiento computacional.} &
\phasecard{3}{milcMagenta}{milcGris}{Praxis\\[-1mm]\small Construyo y aplico.} &
\phasecard{4}{milcVino}{milcGris}{Evaluación\\[-1mm]\small Reflexión liberadora.}
\end{tabularx}


\begin{titlebox}{milcVerde}
Fase 1 · Escuta
\end{titlebox}

\begin{softbox}{milcVerde}{milcGris}{Escena para escuchar}
Camilo abre Excel para registrar los gastos del kiosko de la cafetería. En la primera columna escribe los productos, en la segunda los precios. Pero al sumar, salen errores extraños: aparece \#VALOR! o el resultado es absurdo. Camilo revisa y descubre que escribió \textit{``mil quinientos''} en lugar de 1500. Excel trató eso como texto, no como número.
\end{softbox}

\checkbox Recuerdo una vez que ingresé información en Excel y un cálculo salió mal.\\
\checkbox Identifico la diferencia entre escribir un número y escribirlo como texto.\\
\checkbox Reconozco la importancia de un encabezado claro en cada columna.

\textbf{Una vez ingresé datos en Excel y el cálculo falló porque:}\\[1mm]
\linead{\linewidth}\\[1mm]
\linead{\linewidth}

\textbf{El error fue de tipo o de formato porque:}\\[1mm]
\linead{\linewidth}\\[1mm]
\linead{\linewidth}

\textbf{Lo que aprendí o haría diferente la próxima vez es:}\\[1mm]
\linead{\linewidth}\\[1mm]
\linead{\linewidth}

\begin{infoband}{milcVerde}
\textbf{Dato curioso.} En Excel las fechas son números: el 1 de enero de 1900 es el número 1, el día siguiente es 2, y así. Por eso puedes restar fechas para saber cuántos días pasaron. ``Hoy menos tu cumpleaños'' calcula tu edad en días.
\end{infoband}

\begin{titlebox}{milcTurquesa}
Fase 2 · Sistematización con pensamiento computacional
\end{titlebox}

\begin{softbox}{milcTurquesa}{milcGris}{Cuatro pilares aplicados al tema}
El pensamiento computacional descompone el dato en su tipo, reconoce patrones de formato, abstrae la regla \textit{cada cosa en su columna}, y diseña un algoritmo de ingreso ordenado.
\end{softbox}

\small
\begin{tabularx}{\textwidth}{>{\bfseries\RaggedRight\arraybackslash}p{3.6cm}Y}
\rowcolor{milcTurquesa!90}\color{white}Pilar & \color{white}\bfseries Aplicación al tema\\
1. Descomposición & Descompón cada columna: ¿qué tipo de dato guarda (texto, número, fecha, moneda, porcentaje)?\\
2. Reconocimiento de patrones & Identifica patrones: los números a la derecha, el texto a la izquierda, las fechas con formato uniforme.\\
3. Abstracción & Abstracción: el encabezado y el formato comunican cómo leer la columna sin abrir el dato.\\
4. Algoritmo conceptual & Algoritmo: \textbf{1)} nombrar columnas · \textbf{2)} elegir tipo de cada una · \textbf{3)} aplicar formato · \textbf{4)} ingresar datos · \textbf{5)} verificar con un cálculo de prueba.\\
\end{tabularx}
\normalsize

\begin{drawbox}{milcTurquesa}{Anatomía de una tabla bien formateada}
\textbf{Encabezados} (fila 1): nombres claros, sin abreviaturas raras. \\
\textbf{Tipo por columna}: número, texto, fecha, moneda, porcentaje, lista. \\
\textbf{Formato visual}: alineación (texto izquierda, número derecha). \\
\textbf{Decimales coherentes}: 2 para moneda, 0-1 para edad. \\
\textbf{Verificación}: una fórmula de prueba (=SUMA, =CONTAR) confirma que los tipos están bien.
\end{drawbox}

\begin{softbox}{milcMostaza}{milcGris}{Errores comunes y su corrección}
\begin{itemize}
\item Ingresar números como texto (\textit{``1500''}) $\rightarrow$ Escribe 1500 y aplica formato moneda.
\item Mezclar formatos de fecha (12/03/26 y 2026-03-12) $\rightarrow$ Estandariza ANTES de ingresar.
\item Dejar columnas sin encabezado $\rightarrow$ Cada columna tiene un nombre claro en la fila 1.
\item Centrar todo lo numérico $\rightarrow$ Números a la derecha, texto a la izquierda.
\item No verificar con un cálculo de prueba $\rightarrow$ Una fórmula simple revela errores rápido.
\end{itemize}
\end{softbox}

\textbf{Si fuera a registrar gastos del kiosko, mis 4 columnas serían:}\\[1mm]
\linead{\linewidth}\\[1mm]
\linead{\linewidth}

\textbf{El tipo de dato y formato de cada columna sería:}\\[1mm]
\linead{\linewidth}\\[1mm]
\linead{\linewidth}

\begin{titlebox}{milcMagenta}
Fase 3 · Praxis con pensamiento computacional
\end{titlebox}

\begin{softbox}{milcMagenta}{milcGris}{Construyendo el producto con los cuatro pilares}
Construir una tabla bien formateada es un sistema. Decomponemos los datos, reconocemos patrones de buenas tablas, abstraemos la función de cada columna, y seguimos un algoritmo de armado limpio.
\end{softbox}

\small
\begin{tabularx}{\textwidth}{>{\bfseries\RaggedRight\arraybackslash}p{3.6cm}Y}
\rowcolor{milcMagenta!90}\color{white}Pilar & \color{white}\bfseries Aplicación al producto\\
Descomposición & Tabla: 1) encabezados · 2) columnas con tipo · 3) formato visual coherente · 4) datos reales · 5) verificación con fórmula.\\
Patrones & Patrones de tablas profesionales: encabezados en negrita, alternancia de color en filas, formato uniforme.\\
Abstracción & Función esencial: una tabla bien hecha permite leer e interpretar sin tener que preguntar.\\
Algoritmo de armado & Pasos: definir columnas $\rightarrow$ aplicar formato $\rightarrow$ ingresar 10 registros $\rightarrow$ verificar con SUMA o CONTAR $\rightarrow$ ajustar.\\
\end{tabularx}
\normalsize

\begin{drawbox}{milcMagenta}{Checklist de la tabla bien formateada}
\checkbox 4 columnas mínimo: 1 texto, 1 número, 1 fecha, 1 moneda.\\
\checkbox Encabezados claros en la fila 1.\\
\checkbox Formato aplicado a cada columna (no solo números planos).\\
\checkbox 10 registros reales del entorno escolar.\\
\checkbox Una fórmula de prueba (ej. =SUMA o =CONTAR) verifica los tipos.\\
\checkbox Alineación coherente (texto izquierda, número derecha).
\end{drawbox}

\begin{softbox}{milcMagenta}{milcGris}{Plantilla de trabajo}
\textbf{Encabezados:} Producto · Cantidad · Fecha · Precio. \\
\textbf{Tipo:} Texto · Número · Fecha · Moneda. \\
\textbf{Formato:} Texto izq. · Número der. 0 dec. · Fecha dd/mm/aaaa · Moneda \$ COP 0 dec. \\
\textbf{Datos:} Mínimo 10 productos del kiosko (panes, jugos, frutas, lápices). \\
\textbf{Prueba:} =SUMA(D2:D11) en la fila 12 para verificar moneda.
\end{softbox}

\textbf{Mi tabla del kiosko: 4 columnas con tipos y 10 registros reales.}\\[1mm]
\linead{\linewidth}\\[1mm]
\linead{\linewidth}\\[1mm]
\linead{\linewidth}

\begin{titlebox}{milcMostaza}
Producto de la sesión
\end{titlebox}
\textbf{Tabla con 4 columnas tipificadas + 10 registros del kiosko + fórmula de verificación}.
\begin{softbox}{milcMostaza}{milcGris}{Criterios de calidad}
\begin{itemize}
\item Encabezados claros en fila 1.
\item 4 tipos distintos de dato (texto, número, fecha, moneda).
\item Formato aplicado y verificable.
\item 10 registros reales (no inventados).
\item Una fórmula de prueba que valida los tipos.
\end{itemize}
\end{softbox}

\begin{titlebox}{milcVino}
Fase 4 · Evaluación liberadora — cinco dimensiones
\end{titlebox}

\begin{softbox}{milcVino}{milcGris}{Sentido de la evaluación}
Evaluar no es solo revisar una nota. Es reconocer qué aprendí, qué cambió en mí, qué emociones manejé, cómo me relaciono con la ciudad y la comunidad, y qué le dejaré a quien viene detrás.
\end{softbox}

\textbf{Mi reflexión liberadora en cinco dimensiones.} Completa cada frase iniciada:

\small
\begin{tabularx}{\textwidth}{>{\bfseries\RaggedRight\arraybackslash}p{3.4cm}Y>{\RaggedRight\arraybackslash}p{4.6cm}}
\rowcolor{milcVino}\color{white}Dimensión & \color{white}\bfseries Mi respuesta (frame starter) & \color{white}\bfseries Algo que descubrí\\
1. Personal & Lo que más cambió en mí fue \linead{6.5cm} & \linead{4.2cm}\\
2. Emocional & La emoción más fuerte fue \linead{2.5cm} y la regulé \linead{3cm} & \linead{4.2cm}\\
3. Ciudadana & En democracia esto importa porque \linead{6cm} & \linead{4.2cm}\\
4. Local (Cartago) & En mi barrio sería útil para \linead{6.5cm} & \linead{4.2cm}\\
5. Intergeneracional & A mi abuela/abuelo le diría \linead{6.5cm} & \linead{4.2cm}\\
\end{tabularx}
\normalsize

\begin{infoband}{milcVino}
\textbf{Para la próxima sustentación.} Vuelve a esta tabla en seis meses. Verás que las respuestas cambian — no porque lo aprendido cambie, sino porque tú cambias. La evaluación liberadora es una conversación contigo a través del tiempo.
\end{infoband}


\begin{titlebox}{milcGradeDark}
Triángulo de pensamiento · cierre filosófico
\end{titlebox}

\begin{softbox}{milcVino}{milcGris}{Dussel · lente del nosotros}
\textit{``El Otro es el origen absoluto de la responsabilidad ética.''}

Los datos que registras son personas (sus compras, sus pagos, su tiempo). Una tabla bien hecha respeta a quien aparece en ella: no la confunde, no la deforma, no la expone sin razón.

\textbf{¿De quién son los datos que estás organizando? ¿Pediste permiso si son personales?}
\end{softbox}
\linead{\linewidth}\\[1mm]
\linead{\linewidth}

\begin{softbox}{milcMostaza}{milcGris}{Estoicismo · Marco Aurelio · cuidado interior}
\textit{``El presente es lo único que poseemos. Cuídalo.''}

Cada minuto que pierdes hoy por usar mal los formatos se multiplicará en futuras hojas de cálculo. Diseñar bien una sola vez te ahorra horas mañana. La disciplina del presente cuida tu tiempo futuro.

\textbf{¿Cuántos minutos pierdes hoy por no aplicar formato? ¿Vale la pena hacerlo bien una vez?}
\end{softbox}
\linead{\linewidth}\\[1mm]
\linead{\linewidth}

\begin{softbox}{milcTurquesa}{milcGris}{Floridi · lente de la infoesfera}
\textit{``Toda interacción es un acto de información; no hay neutralidad técnica.''}

Tu elección de \textit{cantidad} o \textit{moneda} en una columna no es estética: comunica algo distinto a quien lee. Decir \textit{``10''} sin formato vs \textit{``\$10.000''} con formato moneda es la diferencia entre confusión y claridad.

\textbf{¿Qué le comunica tu formato a quien va a leer la tabla? ¿Es lo que quieres comunicar?}
\end{softbox}
\linead{\linewidth}\\[1mm]
\linead{\linewidth}

\begin{titlebox}{milcVino}
Compromiso final
\end{titlebox}

\small
\begin{tabularx}{\textwidth}{>{\RaggedRight\arraybackslash}p{3.0cm}YYY}
\rowcolor{milcVino}\color{white}\bfseries Criterio & \color{white}\bfseries Lo logré & \color{white}\bfseries En proceso & \color{white}\bfseries Necesito apoyo\\
Comprensión del tema & Explico ideas centrales con mis palabras. & Reconozco ideas pero falta precisión. & Necesito repasar conceptos base.\\
Pensamiento computacional & Apliqué los 4 pilares al diseñar mi producto. & Apliqué algunos pilares. & Necesito releer la fase 2.\\
Aplicación y producto & Mi producto cumple los criterios y se puede revisar. & Producto funcional con ajustes pendientes. & Aún en construcción.\\
Reflexión 5 dimensiones & Respondí cada dimensión con honestidad. & Respondí algunas con detalle. & Mis respuestas son muy generales.\\
Triángulo de pensamiento & Conecté las 3 voces con mi trabajo real. & Conecté al menos una. & No relaciono las voces con mi proceso.\\
\end{tabularx}
\normalsize

\textbf{Hoy entendí que:}\\[1mm]
\linead{\linewidth}\\[1mm]
\linead{\linewidth}

\textbf{Mi compromiso para la próxima sesión es:}\\[1mm]
\linead{\linewidth}\\[1mm]
\linead{\linewidth}

\begin{softbox}{milcMostaza}{milcGris}{Cita de despedida · Marco Aurelio}
\textit{``El obstáculo en el camino se vuelve el camino. Lo que estorba al camino se vuelve camino.''}\\[1mm]
Cada error de hoy, bien observado, se convierte en pieza del camino que pisarás mañana.
\end{softbox}

\small
\begin{tabularx}{\textwidth}{>{\bfseries}p{3.6cm}Y}
\rowcolor{milcGradeDark!12}Nombre completo: & \linead{10cm}\\
Fecha: & \linead{4cm} \quad Curso/grupo: \linead{4cm}\\
Mi firma: & \linead{9cm}\\
\end{tabularx}
\normalsize

\begin{infoband}{milcGradeDark}
\textbf{Próxima sesión.} Aprenderemos a estructurar tablas más complejas: campos, registros, claves únicas y limpieza de datos. La tabla del tendero se vuelve archivo institucional.
\end{infoband}

\end{document}
